\section{Round 4 (Erd\H{o}s Problem \#10): explicit multiplicity in Crocker's construction}

\subsection{1) ROUND-4 OBJECTIVE}

\noindent\textbf{Path (C): obstruction / correction.}
After Round~3, the main infinite-family obstruction for $k=3$ was rigorous but still only qualitative up to an unspecified implicit constant. The most promising strict advance is therefore to extract \emph{explicit quantitative strength} from Crocker's construction.

\medskip
\noindent
The goal of this round is to prove a theorem strictly stronger than the Round~3 bound $E_3(X)\gg \log\log X$: namely, to show that for each stage of Crocker's construction there are \emph{at least}
\[
C_0:=7827\cdot 2^{554}
\]
distinct even integers, all divisible by $16$, that are not representable as
\[
\text{prime}+\text{at most }3\text{ powers of }2,
\]
and consequently
\[
E_3(X)\ge C_0\bigl(\lfloor \log_2\log_2 X\rfloor-10\bigr)
\qquad (X\ge 2^{2^{11}}).
\]

\subsection{2) Round-3 FOUNDATION USED}

We use the following previously established facts.

\begin{itemize}
\item \textbf{Round~1, Lemma~10.1.} An integer $N$ is $k$-representable if and only if there exists $m\ge 0$ with $\mathrm{wt}_2(m)\le k$ such that $N-m$ is prime.
\item \textbf{Round~1, Lemma~10.2.} If $N$ is even and $N=p+S$ with $p$ an odd prime and $\mathrm{wt}_2(S)\le k$, then $S=1+m$ with $m$ even and $\mathrm{wt}_2(m)\le k-1$, hence $N-1=p+m$.
\item \textbf{Round~2, Lemma~5.1.} If $t\equiv -1\pmod{2^r}$ and $t>2^r-1$, then $\mathrm{wt}_2(t-1)\ge r$; in particular, $t\equiv -1\pmod{16}$ and $t>15$ imply $\mathrm{wt}_2(t-1)\ge 4$.
\item \textbf{Round~3 repaired shift theorem.} If $t$ is one of Crocker's odd integers with
\begin{enumerate}
\item $t\equiv -1\pmod{16}$,
\item $t$ composite,
\item $t\neq p+2^a$ for all primes $p$ and all $a\ge 1$,
\item $t\neq p+2^a+2^b$ for all primes $p$ and all $a,b\ge 1$,
\end{enumerate}
then $N:=t+1$ is even and not $3$-representable.
\end{itemize}

We will \emph{not} re-prove these facts; the new work in this round is entirely quantitative.

\subsection{3) NEW INSIGHT / TOOL (ROUND-4)}

The new ingredient is that Round~3 did not use two explicit numerical outputs of Crocker's proof:

\begin{enumerate}
\item he proves that the parameter $v$ in his counting argument satisfies
\[
(16\prod_{i=1}^{h+1}p_i)\,v<G_k<(v+1)(16\prod_{i=1}^{h+1}p_i),
\]
and then explicitly estimates
\[
16\prod_{i=1}^{h+1}p_i<2^{434}<2^{988}<G_k;
\]
hence
\[
v\ge 2^{554};
\]
\item for the same construction he proves that there are at least
\[
2^{13}-1-13\cdot 28=7827
\]
admissible choices of the residue class $c$.
\end{enumerate}

The key new observation is that \emph{different admissible choices of $c$ produce disjoint congruence classes modulo the full CRT modulus}, so these $7827$ choices multiply the stage-by-stage count.

A second useful observation is that the stage-$s$ odd integers are all divisible by
\[
\frac{2^{2^s}-1}{G_{10}},
\]
hence in particular by the factor $2^{2^{s-1}}+1$ for every $s\ge 11$. This places the resulting even counterexamples inside explicit disjoint tower intervals.

\subsection{4) ATTACK PLAN (ROUND-4)}

Round~3 left two quantitative gaps.

\begin{itemize}
\item First, the constant hidden in $E_3(X)\gg \log\log X$ was not made explicit.
\item Second, the proof counted only one congruence class per stage, whereas Crocker's construction actually allows many admissible choices of $c$.
\end{itemize}

So the plan is:

\begin{enumerate}
\item fix a stage $s\ge 11$ and an admissible choice of $c$;
\item use Crocker's inequalities to show that this single choice contributes at least $v\ge 2^{554}$ odd integers $t$;
\item prove that different admissible $c$ give disjoint sets of $t$, so stage $s$ contributes at least $7827\cdot 2^{554}$ odd integers;
\item invoke the Round~3 shift theorem to convert these into even non-$3$-representable integers;
\item locate them inside the interval
\[
[2^{2^{s-1}}+2,\;2^{2^s}],
\]
and then sum over stages.
\end{enumerate}

This overcomes the previous obstacle because it uses \emph{all} of the multiplicity already present in Crocker's CRT construction, rather than only one residue class per stage.

\subsection{5) WORK (ROUND-4)}

\subsubsection{5.1 Crocker's stage data}

Fix Crocker's explicit choice $k=10$ and his overlapping congruence system (1). For each stage $s>10$, each admissible residue class $c$ modulo
\[
p_{h+1}=2^{13}-1
\]
produces an enlarged simultaneous congruence system $S_s(c)$.
Let
\[
M_s:=16\cdot \frac{2^{2^s}-1}{G_{10}}\cdot \prod_{i=1}^{h+1}p_i.
\]
By the Chinese Remainder Theorem, the solutions of $S_s(c)$ are exactly the integers
\[
t\equiv q_s(c)\pmod{M_s}
\]
for a uniquely determined nonzero residue class $q_s(c)\pmod{M_s}$, with $1\le q_s(c)<M_s$.

Crocker proves that for each such $c$ there are exactly $v$ or $v+1$ positive integers $t\le 2^{2^s}-1$ satisfying $S_s(c)$, where $v$ is independent of $s$ and $c$ and is defined by
\[
(16\prod_{i=1}^{h+1}p_i)\,v<G_{10}<(v+1)(16\prod_{i=1}^{h+1}p_i).
\]
He also proves the explicit estimates
\[
16\prod_{i=1}^{h+1}p_i<2^{434}<2^{988}<G_{10}.
\]
Therefore
\[
\frac{G_{10}}{16\prod_{i=1}^{h+1}p_i}>2^{554}.
\]
Since $v=\left\lfloor \dfrac{G_{10}}{16\prod_{i=1}^{h+1}p_i}\right\rfloor$, it follows that
\[
v\ge 2^{554}.
\]
Hence:

\medskip
\noindent\textbf{Lemma 5.1.}
For every stage $s\ge 11$ and every admissible choice of $c$, the system $S_s(c)$ has at least $2^{554}$ positive solutions $t\le 2^{2^s}-1$.

\textbf{Proof.}
As just noted, each $S_s(c)$ has either $v$ or $v+1$ positive solutions $\le 2^{2^s}-1$, and $v\ge 2^{554}$.
\qed

\subsubsection{5.2 Distinct admissible choices of $c$ give disjoint sets}

Crocker next proves that the number of admissible residue classes $c$ modulo $p_{h+1}=2^{13}-1$ is at least
\[
2^{13}-1-13\cdot 28=7827.
\]
Let $\mathcal C$ be any set of $7827$ such admissible choices.

\medskip
\noindent\textbf{Lemma 5.2.}
Fix $s\ge 11$. If $c_1,c_2\in \mathcal C$ are distinct, then
\[
q_s(c_1)\not\equiv q_s(c_2)\pmod{M_s}.
\]
Consequently, the corresponding solution sets
\[
T_s(c):=\{t\le 2^{2^s}-1:\ t\equiv q_s(c)\pmod{M_s}\}
\]
are pairwise disjoint.

\textbf{Proof.}
In the CRT system $S_s(c)$, the residue class modulo $p_{h+1}$ is exactly $t\equiv c\pmod{p_{h+1}}$.
If $c_1\neq c_2$, then any solution of $S_s(c_1)$ is incongruent modulo $p_{h+1}$ to any solution of $S_s(c_2)$.
Since $p_{h+1}$ divides $M_s$, equality
\[
q_s(c_1)\equiv q_s(c_2)\pmod{M_s}
\]
would force
\[
q_s(c_1)\equiv q_s(c_2)\pmod{p_{h+1}},
\]
hence $c_1\equiv c_2\pmod{p_{h+1}}$, contradiction.
Thus the residue classes modulo $M_s$ are distinct, and therefore the sets $T_s(c)$ are pairwise disjoint.
\qed

\subsubsection{5.3 Explicit stage multiplicity}

Let
\[
T_s:=\bigcup_{c\in\mathcal C} T_s(c).
\]

\textbf{Proposition 5.3.}
For every stage $s\ge 11$,
\[
|T_s|\ge 7827\cdot 2^{554}.
\]

\textbf{Proof.}
By Lemma~5.1, each $T_s(c)$ has cardinality at least $2^{554}$.
By Lemma~5.2, the sets $T_s(c)$ are pairwise disjoint as $c$ ranges over $\mathcal C$.
Hence
\[
|T_s|=\sum_{c\in\mathcal C}|T_s(c)|\ge |\mathcal C|\cdot 2^{554}\ge 7827\cdot 2^{554}.
\]
\qed

\subsubsection{5.4 Conversion to even counterexamples for $k=3$}

\textbf{Theorem 5.4 (uniform stage theorem).}
For every integer $s\ge 11$, the interval
\[
I_s:=[\,2^{2^{s-1}}+2,\;2^{2^s}\,]
\]
contains at least
\[
C_0:=7827\cdot 2^{554}
\]
distinct even integers $N$, all divisible by $16$, that are not representable as a prime plus at most $3$ powers of $2$.

\textbf{Proof.}
Fix $s\ge 11$ and $t\in T_s$.

Because $t$ satisfies $S_s(c)$ for some admissible $c$, the Round~3 repaired shift theorem applies: $t$ is odd, composite, congruent to $-1\pmod{16}$, and is neither of the form $p+2^a$ nor of the form $p+2^a+2^b$ with $a,b\ge 1$. Therefore
\[
N:=t+1
\]
is even and not $3$-representable.

Also $t\equiv -1\pmod{16}$, so $16\mid N$.

It remains to locate $N$ inside $I_s$.
From the defining condition in $S_s(c)$ we have
\[
t\equiv 0\pmod{\frac{2^{2^s}-1}{G_{10}}}.
\]
Now
\[
2^{2^s}-1=\prod_{j=0}^{s-1}(2^{2^j}+1),
\]
and $G_{10}$ is a divisor of the single factor $2^{2^{10}}+1$.
Since $s\ge 11$, the quotient $\dfrac{2^{2^s}-1}{G_{10}}$ still contains the factor $2^{2^{s-1}}+1$.
Hence
\[
t\ge 2^{2^{s-1}}+1.
\]
On the other hand, by construction,
\[
t\le 2^{2^s}-1.
\]
Therefore
\[
2^{2^{s-1}}+2\le N=t+1\le 2^{2^s},
\]
so $N\in I_s$.

Finally, Proposition~5.3 provides at least $C_0$ distinct choices of $t$, hence at least $C_0$ distinct values of $N=t+1$.
\qed

\subsubsection{5.5 Global counting consequence}

Let
\[
E_3(X):=\#\{N\le X:\ N\text{ even and not }3\text{-representable}\}.
\]

\textbf{Corollary 5.5.}
For every $X\ge 2^{2^{11}}$,
\[
E_3(X)\ge C_0\bigl(\lfloor \log_2\log_2 X\rfloor-10\bigr),
\qquad C_0=7827\cdot 2^{554}.
\]
In particular,
\[
E_3(X)\gg \log\log X
\]
with an explicit constant, and indeed the integers counted on the left may all be chosen divisible by $16$.

\textbf{Proof.}
Let
\[
m:=\lfloor \log_2\log_2 X\rfloor.
\]
Then $m\ge 11$ and $2^{2^s}\le X$ for every $11\le s\le m$.
By Theorem~5.4, each interval
\[
I_s=[2^{2^{s-1}}+2,\;2^{2^s}]
\]
contains at least $C_0$ even non-$3$-representable integers.
These intervals are pairwise disjoint.
Hence
\[
E_3(X)\ge \sum_{s=11}^{m} C_0 = C_0(m-10)
= C_0\bigl(\lfloor \log_2\log_2 X\rfloor-10\bigr).
\]
\qed

\subsection{6) ADVERSARIAL VERIFICATION}

\begin{itemize}
\item \textbf{Did we really get a new theorem, rather than a repackaging of Round~3?}
Yes. Round~3 only proved $E_3(X)\gg \log\log X$ with an unspecified constant and no local interval statement. Theorem~5.4 gives an explicit stagewise count $C_0=7827\cdot 2^{554}$ in each tower interval, and Corollary~5.5 gives an explicit global lower bound.

\item \textbf{Is the use of $7827$ legitimate?}
Yes. Crocker explicitly proves that there are at least $7827$ admissible choices for $c$. The only additional point needed here is disjointness of the resulting congruence classes modulo $M_s$, which is established in Lemma~5.2 by reducing modulo $p_{h+1}$.

\item \textbf{Could two different $c$ choices yield the same actual integer $t$ below $2^{2^s}-1$?}
No. Any such $t$ would satisfy both $t\equiv c_1\pmod{p_{h+1}}$ and $t\equiv c_2\pmod{p_{h+1}}$, forcing $c_1=c_2$.

\item \textbf{Is the lower bound $v\ge 2^{554}$ correct?}
Yes. Crocker proves
\[
(16\prod p_i)\,v<G_{10}<(v+1)(16\prod p_i)
\]
and
\[
16\prod p_i<2^{434}<2^{988}<G_{10}.
\]
Therefore
\[
\frac{G_{10}}{16\prod p_i}>2^{554},
\]
so its floor is at least $2^{554}$.

\item \textbf{Why does every stage-$s$ integer lie above $2^{2^{s-1}}+1$?}
Because each stage-$s$ integer is divisible by $\dfrac{2^{2^s}-1}{G_{10}}$, and for $s\ge 11$ this quotient still contains the full factor $2^{2^{s-1}}+1$.

\item \textbf{Have we used any unproved new representation argument?}
No. The passage from odd $t$ to even $N=t+1$ is exactly the repaired Round~3 argument, which itself rests on Round~1 Lemmas~10.1--10.2 and Round~2 Lemma~5.1.
\end{itemize}

\subsection{7) FINAL (EXACTLY ONE)}

\noindent\textbf{UNRESOLVED (BUT STRICTLY ADVANCED).}
This round does not settle Erd\H{o}s Problem~\#10 in full. It does, however, strengthen the current obstruction substantially beyond Round~3:
for every $s\ge 11$, the interval
\[
[2^{2^{s-1}}+2,\;2^{2^s}]
\]
contains at least
\[
7827\cdot 2^{554}
\]
even integers, all divisible by $16$, that are not representable as a prime plus at most $3$ powers of $2$.
Consequently,
\[
E_3(X)\ge 7827\cdot 2^{554}\bigl(\lfloor \log_2\log_2 X\rfloor-10\bigr)
\qquad (X\ge 2^{2^{11}}).
\]
So we now have an explicit, quantitative, infinitely repeated obstruction at $k=3$.

\subsection{8) COMPLETION ESTIMATE (MANDATORY)}

\noindent\textbf{COMPLETION: 64\%}

\subsection{9) REFERENCES}

\begin{enumerate}
\item T. F. Bloom, \emph{Erd\H{o}s Problem \#10}, \texttt{https://www.erdosproblems.com/10}.
\item R. Crocker, \emph{On the sum of a prime and of two powers of two}, Pacific J. Math. \textbf{36} (1971), 103--107.
\end{enumerate}